\scanpage{320}

\mspage{59}
% archive sheet number: 319

choix $\Sigma=\{\pm1\}$ est \deleted{\unsure{assurément|évidemment} \unsure{bien}} à ma connaissance le seul utile, à cause de son lien avec la notion de chemin et le calcul du $\pi_0$ d'un objet de $X$, etc.

7. \uline{Explicitation des catégories} \inserted{\uline{de fractions}} $X\Sigma^{-1}$ en termes \uline{des} $\underline{\mathrm{Ch}}(X;a,b)$.

Si $X$ est une catégorie, on désigne par $\underline{\mathrm{Ch}}(X)$ la catégorie des chemins de $\unsure{x|X}$ (cf \S{} $\deleted{\unsure{5}}\unsure{4|5}$), et pour $a,b\in\operatorname{Ob}X$, par
\[
\underline{\mathrm{Ch}}(X;a,b)\quad\text{ou simplement}\quad\underline{\mathrm{Ch}}(a,b)
\]
celle des chemins de $a$ vers $b$, i.e de source $a$, but $b$. L'\unsure{intention|intuition} géométrique (qu'on explicitera), c'est qu'un tel chemin \inserted{$c$} définit un chemin ``direct'' de $a$ vers $b$ (i.e. un chemin de type $\tau'$ \deleted{\unsure{p|pur}} ``direct'', i.e. \deleted{\unsure{qui}} toutes les flèches de transition sont directes, i.e. $\tau'$ de la forme $\Delta_+^n$), de $=$ longueur que le chemin initial $c$, obtenu en inversant toutes les flèches rétrogrades \unsure{dans} le chemin $c$. Cette opération est effectivement possible dans $X$ lui-même sss la totalité \unsure{flèches|flèche} sont \unsure{inversible|inversibles}, On obtient alors une suite \unsure{composables|composable} de flèches, dont le composé total \deleted{\unsure{est}} (qu'on peut appeler ``l'effet utile'' du \deleted{\unsure{chemin}})
